\documentclass[10pt,a4paper]{amsart}
\usepackage[margin=19mm]{geometry}
\usepackage{amsmath,amssymb,mathtools}
\usepackage[scr=rsfs,cal=euler]{mathalfa}
\usepackage{xfrac}
\usepackage[unicode,hidelinks,bookmarks=false]{hyperref}
\newcommand{\ie}{i.e.\ }
\newcommand{\cf}{cf.\ }
\newcommand{\resp}{respectively\ }
\newcommand{\Subsection}{\subsection}
\providecommand{\nobreakdash}{\nobreak\hbox{-}}
\setlength{\emergencystretch}{2em}
\multlinegap0pt
\usepackage{mathtools}
\usepackage{amssymb}
\usepackage{bm}
\usepackage{xcolor}
\usepackage[shortlabels]{enumitem}
\usepackage{float}
\newtheorem{lemma}{Lemma}
\title{Propagation of chaos for the Landau equation with very soft and Coulomb potentials}
\author{C\^ome Tabary}
\date{Source: arXiv:2506.15795v3 (CC BY 4.0)}
\begin{document}
\maketitle

\noindent\emph{Source and licence.} Propagation of chaos for the Landau equation with very soft and Coulomb potentials. Version arXiv:2506.15795v3 (CC BY 4.0), distributed by arXiv under the Creative Commons Attribution 4.0 International licence, http://creativecommons.org/licenses/by/4.0/, as recorded in the arXiv licence metadata for this exact version and shown on the abstract page https://arxiv.org/abs/2506.15795v3. Full licence notice: \texttt{licenses/arxiv-2506.15795-CC-BY-4.0.txt}.\par\smallskip

\noindent\textit{Editorial scope.} Counted unit: the article's lemma measures (source0.tex lines 1186--1190). The article's complete original proof of it is delivered with it (source0.tex lines 1191--1279, one \texttt{proof} environment of 89 lines). No mathematics is written, repaired or re-derived: the statement and the proof are the article's own lines, in the article's own notation, and no retained line is substituted or re-flowed.  Retained uncounted as the source-local context the proof needs: lines 715--719; lines 761--767.  Nothing was omitted from any retained span, so the package carries no omission record.  No retained line contains a citation, so the wrapper carries no bibliography and no bibliography entry.  No retained line contains a figure environment, an image-inclusion command or drawing code, so no image is delivered and the package declares no asset dependency.  Editorial formatting only, in addition: an AMS \texttt{amsart} preamble that restates the article's own declarations and macros used by the retained lines, namely the environments \texttt{lemma} on separate counters, together with the standard packages those lines need.  The retained environments print in this document as Lemma 1 in order of appearance, whereas the article prints the same items with the numbering of its own full document.  The complete native source of the article as distributed in the arXiv e-print for this exact version is bundled unchanged as \texttt{sources/180109/source0.tex}.
\begin{lemma}
\label{lem:keyestimate}
    For any $N\geq 2$:
    $$\sup_{t\geq 0} \mathbb{E}\left[ \left\vert V^N_1(t) - V^N_2(t)\right\vert^{-2} \right] \leq I_0.$$
\end{lemma}

\begin{lemma}
    \label{lem:holdercont}
    There exists $C_{T,I_0,\gamma}>0$ depending on $T,I_0$ and $\gamma$ such that 
    \begin{align*}
   \mathbb{E}\left[  \sup_{s\neq t\in[0,T]} \frac{d(\mu^N_s,\mu^N_t)}{\vert s-t \vert^\frac{1}{8}} \right]\leq C_{T,I_0,\gamma}.
\end{align*}
\end{lemma}

\begin{lemma}
    \label{aswsolution-empirical measures}
    For any $\varphi\in C^2_b(\mathbb{R}^3)$ and any $t\in[0,T]$,
    $$\mathbb{E}\vert \mathcal{F}_{\varphi,t}(\mu^N)\vert\xrightarrow[N\rightarrow+\infty]{}0.$$
\end{lemma}

\begin{proof}
    \textit{Step 1.} Using that
    $$\int \varphi(v)\mu^N_t(dv)=\frac{1}{N}\sum_{i=1}^N\varphi(V^N_i(t))$$
    by definition of empirical measures, we have:
    \begin{align}
    \label{eq:Fagainstempiricalmeasures1}
    \mathcal{F}_{\varphi,t}&(\mu^N) = -\frac{1}{N} \sum_{i=1}^N\varphi(V^N_i(t)) + \frac{1}{N} \sum_{i=1}^N\varphi(V^N_i(0))\\
    \label{eq:Fagainstempiricalmeasures2}
    &+ \frac{1}{N^2}\int_0^t \sum_{\substack{i,j=1\\ i\neq j}}^N \mathbf{1}_{V^N_i(s)\neq V^N_j(s)}b(V^N_i(s)-V^N_j(s)) \cdot (\nabla\varphi(V^N_i(s))-\nabla\varphi(V^N_j(s)))\ ds\\
    \label{eq:Fagainstempiricalmeasures3}
    &+\frac{1}{N^2} \int_0^t \sum_{\substack{i,j=1\\ i\neq j}}^N \mathbf{1}_{V^N_i(s)\neq V^N_j(s)}\alpha a(V^N_i(s)-V^N_j(s)) \cdot \nabla^2\varphi(V^N_i(s))\ ds.
    \end{align}  
    Note that the cases $i=j$ are excluded by the indicator functions. We now make use that the particle system approximates the Landau equation. Indeed, applying the Itô formula to compute the differential of $t \mapsto \frac{1}{N}\sum_{i=1}^N\varphi(V^N_i(t))$ (as we did at the beginning of the proof of Lemma~\ref{lem:holdercont}), we obtain the following identity, which closely resembles $\mathcal{F}(\mu^N)$:
    \begin{align}
    \label{eq:Itôempiricalmeasures1}
       0 &= -\frac{1}{N}\sum_{i=1}^N\varphi(V^N_i(t))+\frac{1}{N}\sum_{i=1}^N\varphi(V^N_i(0)) \\
       \label{eq:Itôempiricalmeasures2}
       &+ \frac{1}{N(N-1)}\int_0^t \sum_{\substack{i,j=1\\ i\neq j}}^N  b^N\left(V^N_i(s) -V^N_j(s)\right) \cdot (\nabla\varphi\left(V^N_i(s)\right)-\nabla\varphi\left(V^N_j(s)\right)) ds\ \\
       \label{eq:Itôempiricalmeasures3}
       &+ \frac{1}{N(N-1)} \int_0^t \sum_{\substack{i,j=1\\ i\neq j}}^N  \alpha^Na\left(V^N_i(s) -V^N_j(s)\right):\nabla^2\varphi\left(V^N_i(s)\right)ds\\
       \label{eq:Itôempiricalmeasures4}
       &+\sqrt{\frac{2}{N-1}}\frac{1}{N}\sum_{\substack{i,j=1\\ i\neq j}}^N \int_0^t  \nabla \varphi\left(V^N_i(s)\right)\cdot \sigma^N \left(V^N_i(s) - V^N_j(s)\right)dB^N_{ij}(s)
    \end{align}
    We compare each line. Remark that the first lines~\eqref{eq:Fagainstempiricalmeasures1} and~\eqref{eq:Itôempiricalmeasures1} are the same.
    
    \textit{Step 2.} We compare the lines~\eqref{eq:Fagainstempiricalmeasures2} and~\eqref{eq:Itôempiricalmeasures2}, that contain the terms in $b$, which are the most singular. With the shorthand notations $\psi^N_{ij}(s):=\nabla\varphi(V^N_i(s))-\nabla\varphi(V^N_j(s))$, and $\mathbf{1}b(v-w):=\mathbf{1}_{v-w}b(v-w)$ we study the difference
    \begin{align*}
        \frac{1}{N(N-1)}&\int_0^t \sum_{\substack{i,j=1\\ i\neq j}}^N  b^N\left(V^N_i(s) -V^N_j(s)\right) \cdot \psi^N_{ij}(s) ds\\
        &-\frac{1}{N^2}\int_0^t \sum_{\substack{i,j=1\\ i\neq j}}^N \mathbf{1}b(V^N_i(s)-V^N_j(s)) \cdot \psi^N_{ij}(s)\ ds\\=
        &\frac{1}{N(N-1)} \int_0^t \sum_{\substack{i,j=1\\ i\neq j}}^N  (b^N-\mathbf{1}b)\left(V^N_i(s) -V^N_j(s)\right) \cdot \psi^N_{ij}(s) ds \ + \mathcal{R}
    \end{align*}
    with the remainder
    $$\mathcal{R}= \frac{1}{N^2(N-1)} \int_0^T \sum_{\substack{i,j=1\\ i\neq j}}^N \mathbf{1}b(V^N_i(t)-V^N_j(t)) \cdot \psi^N_{ij}(t)\ dt, $$
    which we quickly deal with.
    The estimate $\mathbf{1}b(v)\lesssim \vert v \vert^{\gamma+1}$, the boundedness of $\nabla\varphi$, as well as exchangeability allow us to easily control:
    $$\mathbb{E}\vert \mathcal{R} \vert \leq   \frac{2\Vert \nabla\varphi\Vert_{L^\infty}}{N^2(N-1)} \int_0^t \sum_{\substack{i,j=1\\ i\neq j}}^N \mathbb{E}\vert V^N_i(s)-V^N_j(s) \vert^{\gamma+1}\ ds \leq \frac{2\Vert \nabla\varphi\Vert_{L^\infty}}{N} \int_0^t \mathbb{E}\vert V^N_1(s)-V^N_2(s) \vert^{\gamma+1}\ ds$$
    Using $\vert v \vert^{\gamma+1} \leq 1 + \vert v \vert^{-2}$ and Lemma~\ref{lem:keyestimate} to bound the expectation we get
    $$\mathbb{E}\vert \mathcal{R} \vert \leq \frac{C_{\varphi,I_0,t}}{N}.$$
    Now for the main term, we first write a pointwise estimate, using that $b^N(v)$ and $b(v)$ coincide for $\vert v\vert\geq \eta_N$:
    $$\vert (b^N - \mathbf{1}b)(v)\vert \leq 4 \mathbf{1}_{\vert v \vert \leq \eta_N}\vert v \vert^{\gamma + 1}  \leq 4 \vert v \vert^{-3} \eta_N^{\gamma+4}.$$
    Notice that the above bound also holds for $v=0$ because $b^N(0)=0$, since $b^N(v)=-2\alpha^N(\vert v\vert)v$.
    We rely on the $C^2$ behaviour of $\varphi$ to write
    $$\vert \psi^N_{ij}(s)\vert=\vert\nabla\varphi(V^N_i(s))-\nabla\varphi(V^N_j(s))\vert \leq C_\varphi\vert V^N_i(s)-V^N_j(s)\vert.  $$
    With that estimate we cancel one order of the singularity and we have the following bound:
    \begin{align*}
        \frac{2}{N(N-1)} &\mathbb{E}\left\vert \int_0^t \sum_{\substack{i,j=1\\ i\neq j}}^N  (b^N-\mathbf{1}b)\left(V^N_i(s) -V^N_j(s)\right)  \cdot \psi^N_{ij}(s) ds\right\vert \\
        &\leq \frac{\eta_N^{\gamma+4}C_\varphi}{N(N-1)} \sum_{\substack{i,j=1\\ i\neq j}}^N\int_0^t \mathbb{E} \left[\vert V^N_i(s) -V^N_j(s)\vert^{-2}\right]   ds\\
        &\leq \eta_N^{\gamma+4}C_\varphi \int_0^t \mathbb{E} \left[\vert V^N_1(s) -V^N_2(s)\vert^{-2}\right]   ds\\
        &\leq  \eta_N^{\gamma+4} C_{\varphi,I_0,t}
    \end{align*}
    where the third line is obtained by exchangeability and the last line thanks to Lemma~\ref{lem:keyestimate} to bound the expectation.

    \textit{Step 3.}
    We compare~\eqref{eq:Fagainstempiricalmeasures3} and~\eqref{eq:Itôempiricalmeasures3}, exactly as in the previous step.
    This time we use the two bounds
    $$\Vert\mathbf{1}a(v)\Vert \leq C \vert v \vert^{\gamma+2} \leq C(1 + \vert v \vert^{-2})$$ 
    and
    $$\Vert a^N - \mathbf{1}a(v)\Vert \leq C \mathbf{1}_{\vert v \vert \leq \eta_N}\vert v \vert^{\gamma + 2}  \leq C \vert v \vert^{-2} \eta_N^{\gamma+4}$$
    (which still holds for $v=0$ because $\alpha^N a(0)=0$ since $a(0)=0$.)
    They readily give:
    \begin{align*}
    \mathbb{E}\Bigg\vert &\frac{1}{N^2} \int_0^t \sum_{\substack{i,j=1\\ i\neq j}}^N \mathbf{1}_{V^N_i(s)\neq V^N_j(s)}\alpha a(V^N_i(s)-V^N_j(s)) \cdot \nabla^2_v\varphi(V^N_i(s))\ ds\\
    &- \frac{1}{N(N-1)} \int_0^T \sum_{\substack{i,j=1\\ i\neq j}}^N  \alpha^Na\left(V^N_i(s) -V^N_j(t)\right):\nabla^2_v\varphi\left(V^N_i(s)\right)ds \Bigg\vert\leq C_{\varphi,I_0,t}\left(\eta_N^{\gamma+4}+\frac{1}{N}\right) .
    \end{align*}

    \textit{Step 4.} It remains to control the martingale term~\eqref{eq:Itôempiricalmeasures4}, that we rewrite using $B^N_{ji}=-B^N_{ij}$:
    \begin{align*}
        \mathcal{B} &:= \sqrt{\frac{2}{N-1}}\frac{1}{N}\sum_{\substack{i,j=1\\ i\neq j}}^N \int_0^t  \nabla_v \varphi\left(V^N_i(s)\right)\cdot \sigma^N \left(V^N_i(s) - V^N_j(s)\right)dB^N_{ij}(s)\\
    &= \sqrt{\frac{2}{N^2(N-1)}}\sum_{\substack{i,j=1\\ i< j}}^N \int_0^t  \xi^N_{ij}(s)\cdot \sigma^N \left(V^N_i(s) - V^N_j(s)\right)dB^N_{ij}(s)
    \end{align*}
    since $\sigma^N(-v)=\sigma^N(v)$, with $\psi^N_{ij}(s) = \nabla\varphi(V^N_i(s))-\nabla \varphi(V^N_j(s))$ as before, .
    
    We use independence of the Brownian motions $(B^N_{ij})_{1\leq i<j \leq N}$ to cancel cross terms in the square expectation:
    \begin{align*}
        (\mathbb{E}\vert \mathcal{B}\vert)^2 &\leq\mathbb{E} (\mathcal{B}^2)\\
        &= \frac{2}{N^2(N-1)} \sum_{\substack{i,j=1\\ i< j}}^N \mathbb{E}\left[\int_0^t  \psi^N_{ij}(s)\cdot \sigma^N \left(V^N_i(s) - V^N_j(s)\right)dB^N_{ij}(s) \right]^2
    \end{align*}
    By the Iso isometry we can let the square inside the integral, and we obtain:
    \begin{align*}
        \mathbb{E}\left[\int_0^t  \psi^N_{ij}(s)\cdot \sigma^N \left(V^N_i(s) - V^N_j(s)\right)dB^N_{ij}(s) \right]^2 &\leq C_\varphi \int_0^t\mathbb{E}\left[ \Vert\sigma^N \left(V^N_i(s) - V^N_j(s)\right)\Vert^2 \right]ds
    \end{align*}
    We conclude as before, with the bound $\Vert\sigma^N \left(v\right)\Vert^2 \lesssim \vert v \vert^{\gamma+2} \leq 1 + \vert v \vert^{-2}$, exchangeability and Lemma~\ref{lem:keyestimate} yielding
    \begin{align*}
        (\mathbb{E}\vert \mathcal{B}\vert)^2 &\leq \frac{C_{\varphi,I_0,t}}{N}.
    \end{align*}
    Combining all steps, we have obtained
    $$\mathbb{E}\vert \mathcal{F}_{\varphi,t}(\mu^N)\vert \leq C_{\varphi,I_0,t}\left(\eta_N^{\gamma+3} + \eta_N^{\gamma+2} + \frac{1}{N}+\frac{1}{\sqrt{N}}\right)\xrightarrow[N\rightarrow\infty]{} 0,$$
    and the proof is concluded.
\end{proof}

\end{document}
